%% 01_introduction.tex
%% 序論セクション

\section{はじめに}
\label{sec:introduction}

学術論文の執筆には，書式と構造への細心の注意が必要である．
本テンプレートは，ICSE（International Conference on Software Engineering）投稿用の論文を
ACM sigconf形式で作成するための基盤を提供し，学術コミュニティのベストプラクティスを取り入れている．

\subsection{背景}

ICSEはソフトウェア工学研究における主要な発表の場である．
適切に設計されたテンプレートにより，著者はACM出版基準に準拠した一貫した書式を確保しながら
内容に集中することができる．

\subsection{機能}

本テンプレートには，以下の便利な機能が含まれている：

\begin{itemize}
    \item ACM sigconf形式に基づく2カラムレイアウト
    \item 定理環境を含む数式組版
    \item シンタックスハイライト付きコードリスト
    \item アルゴリズム記述機能
    \item ACM Reference Format による自動参考文献管理
\end{itemize}

\subsection{数式の例}

本テンプレートは様々な数学的構造をサポートしている．
例えば，簡単な式：

\begin{equation}
    E = mc^2
    \label{eq:einstein}
\end{equation}

より複雑な整列式：

\begin{align}
    \nabla \cdot \mathbf{E} &= \frac{\rho}{\varepsilon_0} \\
    \nabla \cdot \mathbf{B} &= 0 \\
    \nabla \times \mathbf{E} &= -\frac{\partial \mathbf{B}}{\partial t} \\
    \nabla \times \mathbf{B} &= \mu_0 \mathbf{J} + \mu_0 \varepsilon_0 \frac{\partial \mathbf{E}}{\partial t}
\end{align}

\begin{theorem}
\label{thm:example}
任意の実数 $a$ と $b$ に対して，$(a + b)^2 = a^2 + 2ab + b^2$ が成り立つ．
\end{theorem}

\begin{proof}
$(a + b)(a + b)$ を直接展開することで得られる．
\end{proof}

\subsection{コードリストの例}

本テンプレートはシンタックスハイライト付きのコードリストをサポートしている：

\begin{lstlisting}[language=Python, caption={Pythonコードの例}]
def hello_world():
    """挨拶メッセージを出力する"""
    print("Hello, World!")
    return True

if __name__ == "__main__":
    hello_world()
\end{lstlisting}

\subsection{アルゴリズムの例}

algorithm2e環境を使用してアルゴリズムを記述できる：

\begin{algorithm}[t]
\caption{アルゴリズムの例}
\label{alg:example}
\KwIn{入力データ $X$}
\KwOut{処理結果 $Y$}
$Y \leftarrow \emptyset$ で初期化\;
\For{$X$ の各要素 $x$ について}{
    $x$ を処理\;
    結果を $Y$ に追加\;
}
\Return{$Y$}
\end{algorithm}

\subsection{論文の構成}

本論文の残りの構成は以下の通りである．
\ref{sec:conclusion}節では，本研究の貢献のまとめと
今後の研究の方向性について述べる．
